\documentclass[11pt,a4paper]{article}

%% Packages
\usepackage[landscape, margin=1.8cm, top=2cm, bottom=2cm]{geometry}
\usepackage[T1]{fontenc}
\usepackage[utf8]{inputenc}
\usepackage[activate=false]{microtype}
\usepackage{xcolor}
\usepackage{colortbl}
\usepackage{booktabs}
\usepackage{longtable}
\usepackage{array}
\usepackage{ragged2e}
\usepackage{multirow}
\usepackage{setspace}
\usepackage{parskip}
\usepackage{titlesec}
\usepackage{tcolorbox}
\usepackage[
  colorlinks=true,
  linkcolor=NavyBlue,
  urlcolor=NavyBlue,
  citecolor=NavyBlue,
  pdftitle={Biological Validation of Top Pathway Importance Findings},
]{hyperref}

%% ── Colour-blind-friendly palette (Paul Tol "bright" + IBM accessible) ───────
%% All hues distinguishable under deuteranopia, protanopia, tritanopia.

\definecolor{NavyBlue}    {HTML}{1B3A6B}
\definecolor{MidBlue}     {HTML}{2B5BA8}
\definecolor{BorderGrey}  {HTML}{BBBBBB}
\definecolor{SubtleGrey}  {HTML}{666666}
\definecolor{White}       {HTML}{FFFFFF}
\definecolor{TextDark}    {HTML}{222222}
\definecolor{AccentVerd}  {HTML}{117733}   % dark green for "Supported" text
\definecolor{HeaderBg}    {HTML}{333333}   % near-black main header

%% TMT group — blue family  (#4477AA Tol bright-blue, lightened)
\definecolor{TMTheadBg}   {HTML}{4477AA}   % header band
\definecolor{TMTrow}      {HTML}{D6E4F2}   % row fill
\definecolor{TMTquote}    {HTML}{BDD4E8}   % quote cell (slightly deeper)

%% RT group  — orange family (#EE7733 Tol bright-orange, lightened)
\definecolor{RTheadBg}    {HTML}{CC6600}   % header band
\definecolor{RTrow}       {HTML}{FAE5C8}   % row fill
\definecolor{RTquote}     {HTML}{F5D0A0}   % quote cell

%% OS group  — green family (#228833 Tol bright-green, lightened)
\definecolor{OSheadBg}    {HTML}{228833}   % header band
\definecolor{OSrow}       {HTML}{D2EDD7}   % row fill
\definecolor{OSquote}     {HTML}{B3DDB9}   % quote cell

%% ── Section styling ──────────────────────────────────────────────────────────
\titleformat{\section}{\large\bfseries\color{NavyBlue}}{}{0em}{}[%
  \vspace{-4pt}\textcolor{MidBlue}{\rule{\linewidth}{0.6pt}}]

%% ── tcolorbox ────────────────────────────────────────────────────────────────
\tcbuselibrary{skins,breakable}

%% ── Column types ─────────────────────────────────────────────────────────────
\newcolumntype{N}{>{\centering\arraybackslash\small}p{0.45cm}}
\newcolumntype{P}{>{\centering\arraybackslash\small}p{2.3cm}}
\newcolumntype{C}{>{\centering\arraybackslash\small}p{1.3cm}}
\newcolumntype{H}{>{\centering\arraybackslash\small}p{1.6cm}}
\newcolumntype{Q}{>{\RaggedRight\arraybackslash\footnotesize\itshape}p{9.2cm}}
\newcolumntype{L}{>{\RaggedRight\arraybackslash\footnotesize}p{8.8cm}}
\newcolumntype{V}{>{\centering\arraybackslash\small\bfseries\color{AccentVerd}}p{1.7cm}}

%% ── Group separator row ──────────────────────────────────────────────────────
\newcommand{\grouprow}[2]{%
  \multicolumn{7}{>{\columncolor{#1}}l}{%
    \hspace{6pt}\textbf{\color{White}\small #2}%
  }\\[-2pt]
}

\begin{document}
\setstretch{1.15}

%% ── Title ────────────────────────────────────────────────────────────────────
\begin{center}
{\LARGE\bfseries\color{NavyBlue}
Biological Validation of Pathway Importance Findings from\\[4pt]
Pathway-Informed Deep Learning Models Across Solid Tumour Cohorts}\\[7pt]
{\large\color{MidBlue}
Supplementary Evidence Table~S1 --- TherapAgent Benchmark Study (ACM-BCB \textquotesingle26)}\\[5pt]
{\normalsize\color{SubtleGrey}
TCGA Pan-Cancer Cohorts ($n = 2{,}622$)
$\;\cdot\;$ Reactome Pathway Analysis (1,706 pathways)
$\;\cdot\;$ BINN \& GraphPath Architectures}\\[5pt]
\textcolor{BorderGrey}{\rule{\linewidth}{1.2pt}}
\end{center}

\vspace{4pt}

%% ── Colour legend ─────────────────────────────────────────────────────────────
\begin{center}
\small
\colorbox{TMTheadBg}{\color{White}\textbf{\ \ TMT\ \ }}\ Targeted Molecular Therapy
\qquad
\colorbox{RTheadBg}{\color{White}\textbf{\ \ RT\ \ \ }}\ Radiation Therapy
\qquad
\colorbox{OSheadBg}{\color{White}\textbf{\ \ OS\ \ \ }}\ Overall Survival ($\geq$\,180 days)
\\[3pt]
{\footnotesize\color{SubtleGrey}
Colours follow the Paul Tol colour-blind-safe palette (blue / orange / green),
distinguishable under deuteranopia, protanopia, and tritanopia.}
\end{center}

\vspace{6pt}

%% ── Table ────────────────────────────────────────────────────────────────────
\section*{Table~S1.\enspace Biological Validation of Top-Ranked Pathway Importance Scores
  by Clinical Prediction Head and Cancer Cohort}

\begingroup
\small
\setlength{\arrayrulewidth}{0.4pt}
\arrayrulecolor{BorderGrey}
\renewcommand{\arraystretch}{1.65}

\begin{longtable}{N P C H Q L V}

%% ── Column header (first page) ───────────────────────────────────────────────
\rowcolor{HeaderBg}
\multicolumn{1}{>{\centering\arraybackslash\bfseries\color{White}\small}p{0.45cm}}{\#} &
\multicolumn{1}{>{\centering\arraybackslash\bfseries\color{White}\small}p{2.3cm}}{Pathway} &
\multicolumn{1}{>{\centering\arraybackslash\bfseries\color{White}\small}p{1.3cm}}{Clinical Head} &
\multicolumn{1}{>{\centering\arraybackslash\bfseries\color{White}\small}p{1.6cm}}{Cancer Cohort} &
\multicolumn{1}{>{\centering\arraybackslash\bfseries\color{White}\small}p{9.2cm}}{Evidence Quote} &
\multicolumn{1}{>{\centering\arraybackslash\bfseries\color{White}\small}p{8.8cm}}{Paper Reference and Link} &
\multicolumn{1}{>{\centering\arraybackslash\bfseries\color{White}\small}p{1.7cm}}{Verdict}
\\[2pt]
\midrule
\endfirsthead

\multicolumn{7}{l}{\small\itshape\color{SubtleGrey}%
  Table~S1 continued --- Biological Validation of Top Pathway Importance Findings}\\[2pt]
\rowcolor{HeaderBg}
\multicolumn{1}{>{\centering\arraybackslash\bfseries\color{White}\small}p{0.45cm}}{\#} &
\multicolumn{1}{>{\centering\arraybackslash\bfseries\color{White}\small}p{2.3cm}}{Pathway} &
\multicolumn{1}{>{\centering\arraybackslash\bfseries\color{White}\small}p{1.3cm}}{Clinical Head} &
\multicolumn{1}{>{\centering\arraybackslash\bfseries\color{White}\small}p{1.6cm}}{Cancer Cohort} &
\multicolumn{1}{>{\centering\arraybackslash\bfseries\color{White}\small}p{9.2cm}}{Evidence Quote} &
\multicolumn{1}{>{\centering\arraybackslash\bfseries\color{White}\small}p{8.8cm}}{Paper Reference and Link} &
\multicolumn{1}{>{\centering\arraybackslash\bfseries\color{White}\small}p{1.7cm}}{Verdict}
\\[2pt]
\midrule
\endhead

\midrule
\multicolumn{7}{r}{\small\itshape\color{SubtleGrey}Continued on next page\ldots}
\endfoot

\bottomrule
\endlastfoot

%% ═══════════════════════════════════════════════════════════════════
%%  GROUP 1 — TMT  (blue)
%% ═══════════════════════════════════════════════════════════════════
\grouprow{TMTheadBg}{Group 1 \textbar{} Targeted Molecular Therapy (TMT)}\\[1pt]

\rowcolor{TMTrow}
1 &
ERK/MAPK Targets &
TMT &
Lung &
\cellcolor{TMTquote}``The MAPK pathway (Ras, Raf, and MEK) has recently been discovered to be
a feasible target for innovative cancer therapy\ldots\ Targeting the Ras/Raf/MEK/ERK pathway
is a promising and alternative method in NSCLC treatment.'' &
\textbf{Everything Old Is New Again: Drug Repurposing Approach for NSCLC Targeting MAPK Signaling Pathway}
\newline
Sisinthy \& Bhatt. \textit{Frontiers in Oncology}, 2021.
\newline
\url{https://doi.org/10.3389/fonc.2021.741326} &
\textbf{Supported} \\[3pt]

%% ═══════════════════════════════════════════════════════════════════
%%  GROUP 2 — RT  (orange)
%% ═══════════════════════════════════════════════════════════════════
\grouprow{RTheadBg}{Group 2 \textbar{} Radiation Therapy (RT)}\\[1pt]

\rowcolor{RTrow}
2 &
Receptor-Mediated Mitophagy &
RT &
Breast &
\cellcolor{RTquote}``Mitophagy had an important role in maintaining hypoxia-induced
radioresistance\ldots\ By interacting with Parkin, p53 inhibited the translocation of Parkin
to the mitochondria, disrupting the protective mitophagy process.'' &
\textbf{TAT-ODD-p53 Enhances Radiosensitivity of Hypoxic Breast Cancer Cells by Inhibiting Parkin-Mediated Mitophagy}
\newline
Zheng et al. \textit{Oncotarget}, 2015.
\newline
\url{https://www.ncbi.nlm.nih.gov/pmc/articles/PMC4627318/} &
\textbf{Supported} \\[3pt]

\rowcolor{RTrow}
3 &
ERK/MAPK Targets &
RT &
Head \& Neck &
\cellcolor{RTquote}``We and others have shown that the MAPK/ERK axis plays a significant
role in the development of radioresistance in HNSCC. IR mediates pathway activation, which
leads to reduced radiosensitivity.'' &
\textbf{Adaptive ERK Signalling Activation in Response to Therapy and In Silico Prognostic Evaluation of EGFR-MAPK in HNSCC}
\newline
Binder et al. \textit{British Journal of Cancer}, 2020.
\newline
\url{https://doi.org/10.1038/s41416-020-0892-9} &
\textbf{Supported} \\[3pt]

\rowcolor{RTrow}
4 &
VEGFR2-Mediated Vascular Permeability &
RT &
Lung &
\cellcolor{RTquote}``Molecular inhibition of VEGFR2 could enhance tumor radiation response
through molecular targeting of tumor vasculature. Paracrine signaling from host VEGF to
cancer cell VEGFR2 might be a significant component of RT failures.'' &
\textbf{MiR-200c Increases the Radiosensitivity of NSCLC Cell Line A549 by Targeting VEGF-VEGFR2 Pathway}
\newline
Ge et al. \textit{PLoS ONE}, 2013.
\newline
\url{https://www.ncbi.nlm.nih.gov/pmc/articles/PMC3813610/} &
\textbf{Supported} \\[3pt]

\rowcolor{RTrow}
5 &
Receptor-Mediated Mitophagy &
RT &
Prostate &
\cellcolor{RTquote}``Autophagy contributes to radioresistance by limiting the generation
of reactive oxygen species (ROS) from dysfunctional mitochondria, clearing radiation-induced
damaged macromolecules, and recycling cellular components to sustain the ATP and nucleotide
pools required for DNA repair.'' &
\textbf{Androgen Receptor Contributes to Radioresistance Through DNA Repair and Autophagy in AR-Positive Prostate Cancer Cells}
\newline
Kudo et al. \textit{bioRxiv}, 2025.
\newline
\url{https://doi.org/10.64898/2025.12.03.690226} &
\textbf{Supported} \\[3pt]

%% ═══════════════════════════════════════════════════════════════════
%%  GROUP 3 — OS  (green)
%% ═══════════════════════════════════════════════════════════════════
\grouprow{OSheadBg}{Group 3 \textbar{} Overall Survival $\geq$ 6 Months (OS)}\\[1pt]

\rowcolor{OSrow}
6 &
ERK/MAPK Targets &
OS &
Lung &
\cellcolor{OSquote}``p-ERK1/2-positive expression was an independent prognostic factor for
poor overall survival (OS) in NSCLC patients on both univariate analysis.'' &
\textbf{Prognostic Values of ERK1/2 and p-ERK1/2 Expressions for Poor Survival in Non-Small Cell Lung Cancer}
\newline
Zhao et al. \textit{Tumor Biology}, Springer, 2015.
\newline
\url{https://doi.org/10.1007/s13277-015-3048-4} &
\textbf{Supported} \\[3pt]

\rowcolor{OSrow}
7 &
ERK/MAPK Targets &
OS &
Head \& Neck &
\cellcolor{OSquote}``HNSCC patients with high intratumor expressions of p-MAPK1/3 (p-ERK1/2)
have poor survival, and inhibition of p-ERK1/2 by MAPK pathway inhibitors often inhibited
HNSCC cell growth in vitro.'' &
\textbf{Precision Drugging of the MAPK Pathway in Head and Neck Cancer}
\newline
Fisch et al. \textit{npj Genomic Medicine}, Nature, 2022.
\newline
\url{https://doi.org/10.1038/s41525-022-00293-1} &
\textbf{Supported} \\[3pt]

\end{longtable}
\endgroup

\vspace{4pt}

%% ── Abbreviations ────────────────────────────────────────────────────────────
\begin{tcolorbox}[
  enhanced,
  colback=white,
  colframe=BorderGrey,
  boxrule=0.5pt,
  arc=3pt,
  left=6pt,right=6pt,top=5pt,bottom=5pt,
]
\footnotesize
\textbf{Abbreviations.}
\textbf{TMT} = Targeted Molecular Therapy;\quad
\textbf{RT} = Radiation Therapy;\quad
\textbf{OS} = Overall Survival ($\geq$\,180 days);\quad
\textbf{HNSCC} = Head and Neck Squamous Cell Carcinoma;\quad
\textbf{NSCLC} = Non-Small Cell Lung Cancer;\quad
\textbf{MAPK} = Mitogen-Activated Protein Kinase;\quad
\textbf{VEGFR2} = Vascular Endothelial Growth Factor Receptor 2;\quad
\textbf{TCGA} = The Cancer Genome Atlas;\quad
\textbf{ssGSEA} = Single-Sample Gene Set Enrichment Analysis.\\[3pt]
Pathway importance scores were derived from normalised attribution weights produced by the BINN
and GraphPath architectures following the evaluation protocol described in the manuscript.
Evidence quotes are verbatim extracts from the cited publications.
All DOI and PMC links are clickable in a compatible PDF viewer.
\end{tcolorbox}

\end{document}
